% !TEX root = main.tex
\documentclass[main.tex]{subfiles}
\begin{document}

\appendix

\section{Derivatives of $\energy$ and $(\dual,\combweight)$}
\label{sec:appendix-prelim}

Both our objectives $\centroiddual$ and $\jointdual$ are the plain power diagram energy $\energy$ in \eqref{eq:power-diagram} evaluated at the shifted weighted sites $(\dual, \combweight)$ of \eqref{eq:shifted-site}.
We summarize here the ingredients shared by the next two sections: the Jacobian of \eqref{eq:shifted-site} and the derivatives of \eqref{eq:power-diagram}.

The first-order Jacobian of the reparametrization \eqref{eq:shifted-site} are
\begin{equation}
  \partial_{\shift_i} \dual_i = \Id , \qquad
  \partial_{\shift_i} \combweight_i = 2\,(\dual_i - \tcentroid_i) , \qquad
  \partial_{\weight_i} \combweight_i = 1 ,
  \label{eq:appendix-jac}
\end{equation}
and its only nonzero second-order part is $\smash{\partial^2_{\shift_i} \combweight_i = 2\,\Id}$.

As detailed in \cite[\S2.4--2.5]{deGoes2014Thesis}, the first derivatives of the energy $\energy$ in \eqref{eq:power-diagram} are the standard weight and site gradients:
\begin{equation}
  \partial_{\dual_i} \energy = 2 \int_{\cell_i} (\dual_i - \xx) \, d\xx ,
  \qquad
  \partial_{\combweight_i} \energy = -\cellvol_i ,
  \label{eq:appendix-energy-deriv}
\end{equation}
and its second derivatives are facet integrals with off-diagonals
\begin{subequations}
  \label{eq:appendix-energy-hess-off}
  \begin{empheq}[left={\empheqlbrace\,}]{align}
    \partial^2_{\dual_i \dual_j} \energy
      &= \frac{2}{\norm{\dual_i - \dual_j}} \int_{\face_{ij}}\! (\xx - \dual_i)(\xx - \dual_j)\tr \, dA , \label{eq:appendix-energy-hess-pp} \\
    \partial^2_{\dual_i \combweight_j} \energy
      &= \frac{1}{\norm{\dual_i - \dual_j}} \int_{\face_{ij}}\! (\xx - \dual_i) \, dA , \label{eq:appendix-energy-hess-pl} \\
    \partial^2_{\combweight_i \combweight_j} \energy
      &= \frac{1}{2}\frac{\area_{ij}}{\norm{\dual_i - \dual_j}} , \label{eq:appendix-energy-hess-ll}
  \end{empheq}
\end{subequations}
and diagonals
\begin{subequations}
  \label{eq:appendix-energy-hess-diag}
  \begin{empheq}[left={\empheqlbrace\,}]{align}
    \partial^2_{\dual_i \dual_i} \energy
      &= 2\,\cellvol_i\,\Id - \!\sum_{j \sim i} \frac{2}{\norm{\dual_i - \dual_j}} \int_{\face_{ij}}\! (\xx - \dual_i)(\xx - \dual_i)\tr \, dA , \label{eq:appendix-energy-hess-pp-diag} \\
    \partial^2_{\dual_i \combweight_i} \energy
      &= -\!\sum_{j \sim i} \frac{1}{\norm{\dual_i - \dual_j}} \int_{\face_{ij}}\! (\xx - \dual_i) \, dA , \label{eq:appendix-energy-hess-pl-diag} \\
    \partial^2_{\combweight_i \combweight_i} \energy
      &= -\!\sum_{j \sim i} \frac{1}{2}\frac{\area_{ij}}{\norm{\dual_i - \dual_j}} , \label{eq:appendix-energy-hess-ll-diag}
  \end{empheq}
\end{subequations}
where the site--site diagonal carries an extra interior term $2\,\cellvol_i\,\Id$ from the explicit dependence of $\partial_{\dual_i}\energy$ on $\dual_i$, and otherwise each diagonal is the negated neighbor sum of its off-diagonal counterpart.

\section{Derivatives and concavity of $\centroiddual$}
\label{sec:appendix-centroid}

For fixed weighted sites $(\qq, \weight)$, the energy $\centroiddual(\shift) = \energy(\dual, \combweight)$ in \eqref{eq:centroid-dual} depends on the shifts $\shift$ \emph{only} through the shifted sites $(\dual, \combweight)$ of \eqref{eq:shifted-site}.
Applying the chain rule with the Jacobian \eqref{eq:appendix-jac} and the energy gradients \eqref{eq:appendix-energy-deriv}, we obtain the first derivative of $\centroiddual$:
\begin{equation}
  \begin{aligned}
    \partial_{\shift_i} \centroiddual
      &= (\partial_{\shift_i} \dual_i) \, (\partial_{\dual_i} \energy)
         + (\partial_{\shift_i} \combweight_i) \, (\partial_{\combweight_i} \energy) \\
      &= (\partial_{\dual_i} \energy) + 2\,(\dual_i - \tcentroid_i)\, (\partial_{\combweight_i} \energy)
       = 2 \int_{\cell_i} (\tcentroid_i - \xx) \, d\xx ,
  \end{aligned}
  \label{eq:appendix-centroid-grad}
\end{equation}
thus recovering the gradient \eqref{eq:centroid-dual-gradient}.
% 
For its Hessian, we differentiate \eqref{eq:appendix-centroid-grad} once more, again via the chain rule.
Letting $\alpha, \beta$ range over the shifted-site variables $\{\dual, \combweight\}$, this yields
\begin{equation}
  \partial^2_{\shift_i \shift_j} \centroiddual
    = \sum\nolimits_{\alpha, \beta} (\partial_{\shift_i} \alpha_i)\,(\partial^2_{\alpha_i \beta_j} \energy)\,(\partial_{\shift_j} \beta_j)
      + \delta_{ij}\,(\partial^2_{\shift_i} \combweight_i)\,(\partial_{\combweight_i} \energy) ,
  \label{eq:appendix-centroid-hess-chain}
\end{equation}
with $\delta_{ij}$ indicating the Kronecker delta.
Substituting the energy Hessian \eqref{eq:appendix-energy-hess-off}--\eqref{eq:appendix-energy-hess-diag} and the Jacobian \eqref{eq:appendix-jac}, 
% and using $\smash{\xx - \tcentroid_i = (\xx - \dual_i) + (\dual_i - \tcentroid_i)}$, 
the facet terms combine into $\smash{(\xx - \tcentroid_i)(\xx - \tcentroid_j)\tr}$, recovering the facet Laplacian in \eqref{eq:centroid-hessian}.

% \paragraph{Negative semidefiniteness and null-space.}
To read off the definiteness of \eqref{eq:centroid-hessian}, we introduce the auxiliary per-cell linear field $\varphi_i(\xx) = 2\,\shift_i\tr (\xx - \tcentroid_i)$, whose variation under a shift perturbation $\delta\shift_i$ is $\delta\varphi_i(\xx) = 2\,\delta\shift_i\tr (\xx - \tcentroid_i)$.
Contracting the Hessian blocks with a perturbation $\smash{\delta\shift = \{\delta\shift_i\}}$ collapses each facet to a scalar square:
\begin{equation*}
  \delta\shift\tr \nabla^2 \centroiddual \, \delta\shift
  \!=\! 
  -\sum\nolimits_{i,j}\! \frac{1}{2\,\norm{\dual_i \!-\! \dual_j}} \int_{\face_{ij}}\! \bigl( \delta\varphi_i(\xx) - \delta\varphi_j(\xx) \bigr)^2 \, dA
  \le 0 ,
  \label{eq:appendix-centroid-nsd}
\end{equation*}
so $\nabla^2 \centroiddual \preceq 0$ and $\centroiddual$ is concave.
Being a sum of squares, this quadratic form vanishes only when $\delta\varphi_i(\xx) = \delta\varphi_j(\xx)$ on every shared facet $\face_{ij}$.
Note that this condition does \emph{not} glue the affine fields $\delta\varphi_i$ into a single global affine function: two affine functions that agree on the $(d\!-\!1)$-dimensional facet $\face_{ij}$ only agree on its supporting hyperplane, and their difference may be any multiple of the affine function defining this hyperplane.
To characterize the null-space, we name $a_i$ the affine function of the lifting in \S\ref{sec:background}, evaluated at the shifted weighted site of \eqref{eq:shifted-site}:
\begin{equation}
  a_i(\xx) = 2\,\dual_i\tr\xx - \norm{\dual_i}^2 + \combweight_i = 2\,\qq_i\tr\xx - \norm{\qq_i}^2 + 2\,\shift_i\tr(\xx - \tcentroid_i) + \weight_i ,
  \label{eq:appendix-affine}
\end{equation}
so the facet $\face_{ij}$ lies on the hyperplane $\{a_i = a_j\}$ and a perturbation induces $\delta a_i = \delta\varphi_i$ for fixed weights.
The null-space condition then reads $\smash{\delta a_i - \delta a_j = \kappa_{ij}\,(a_i - a_j)}$ for scalars $\kappa_{ij}$, \ie the perturbation keeps every facet hyperplane fixed to first order and, therefore, leaves the cells unchanged to first order.

In terms of the lifting of \S\ref{sec:background}, we have $\smash{a_i(\xx) = 2\,\dual_i\tr\xx - \sheight_i}$, where the height $\smash{\sheight_i = \norm{\dual_i}^2 - \combweight_i}$ of the shifted weighted site is related to the height $\smash{\height_i = \norm{\qq_i}^2 - \weight_i}$ of the original one by $\smash{\sheight_i = \height_i + 2\,\tcentroid_i\tr\shift_i}$.
Each shift thus moves the lifted point $\smash{(\dual_i, \sheight_i) = (\qq_i, \height_i) + (\shift_i, 2\,\tcentroid_i\tr\shift_i)}$ affinely, and \eqref{eq:appendix-affine} gives $\smash{a_i(\tcentroid_i) = 2\,\qq_i\tr\tcentroid_i - \height_i = \gamma_i}$ for every shift, with $\gamma_i$ as in \eqref{eq:pivot}.
In other words, the shift $\shift_i$ pivots the lifted hyperplane of site $i$ about the fixed point $\smash{(\tcentroid_i, \gamma_i)}$.
Applying the homotheties of \eqref{eq:lifted-homothety} to the lifted shifted sites and differentiating at $s \!=\! 1 \!+\! \mu$ with an infinitesimal offset $(\translation, -c)$ gives the infinitesimal homotheties
\begin{equation}
  \delta\dual_i = \mu\,\dual_i + \translation ,
  \qquad
  \delta\sheight_i = \mu\,\sheight_i - c ,
  \label{eq:infinitesimal-homothety}
\end{equation}
\ie $\smash{\delta a_i = \mu\,a_i + 2\,\translation\tr\xx + c}$, which always satisfy the null-space condition with $\kappa_{ij} \equiv \mu$.
As is generic \cite{Bourne2025Inverting}, we assume that they are also the only solutions, \ie that the cells determine their lifted affine functions only up to these homotheties; this fails in special configurations, such as strips of parallel cells or a boundary cell with a single neighbor, whose affine function may rotate about their shared facet.

With fixed weights, a shift perturbation changes each affine function by $\smash{\delta a_i = 2\,\delta\shift_i\tr(\xx - \tcentroid_i)}$, which vanishes at $\tcentroid_i$, so an infinitesimal homothety is realizable by the shifts exactly when it preserves all pivots.
Matching the linear and constant parts of $\delta a_i = 2\,\delta\shift_i\tr(\xx - \tcentroid_i)$ with $\mu\,a_i + 2\,\translation\tr\xx + c$ gives
\begin{equation}
  \delta\shift_i = \mu\,\dual_i + \translation ,
  \qquad
  0 = \mu\,\gamma_i + 2\,\translation\tr\tcentroid_i + c ,
  \label{eq:appendix-centroid-kernel}
\end{equation}
where the second identity evaluates the first at $\smash{\xx = \tcentroid_i}$.
Since $\delta\shift = 0$ forces $\mu = 0$ and $\translation = 0$ when the shifted sites are distinct, and then $c = 0$, the Hessian is \emph{negative definite} if and only if \eqref{eq:appendix-centroid-kernel} has only the trivial solution, \ie if and only if the vectors $\smash{(\gamma_i, \tcentroid_i, 1)}$ span $\R^{d+2}$.
Equivalently, the pivots $\smash{(\tcentroid_i, \gamma_i) \in \R^{d+1}}$ do not lie on a common hyperplane, as stated in \S\ref{sec:centroid-constraints}.
Since $n$ points in $\R^{d+1}$ always lie on a common hyperplane when $n \!\le\! d\!+\!1$, this requires $n \!\ge\! d\!+\!2$.
A vertical hyperplane ($\mu \!=\! 0$) corresponds to targets $\tcentroid_i$ lying on a common hyperplane of $\R^d$.
In the centroidal setting $\qq_i \!=\! \tcentroid_i$ with $\weight_i \!=\! 0$, the pivots $\smash{(\tcentroid_i, \norm{\tcentroid_i}^2)}$ lie on the paraboloid, and a non-vertical hyperplane ($\mu \!\neq\! 0$) cuts it along a sphere, recovering the Delaunay degeneracy of targets on a common hyperplane or sphere.
In general, $\gamma_i$ lowers the paraboloid by the power distance from $\tcentroid_i$ to its own weighted site, so the degeneracy is that of the regular triangulation of the targets weighted by these power distances.
In particular, the condition that the targets affinely span $\R^d$ is not sufficient.

\section{Derivatives and concavity of $\jointdual$}
\label{sec:appendix-joint}

The joint dual $\jointdual(\shift, \weight) = \energy(\dual, \combweight) + \sum\nolimits_i \weight_i \pvol_i$ defined in \eqref{eq:joint-dual} again depends on $(\shift, \weight)$ only through the shifted sites \eqref{eq:shifted-site}, plus the explicit term $\weight_i \pvol_i$.
Relative to \S\ref{sec:appendix-centroid}, the only new ingredient is the weight reparametrization $\partial_{\weight_i} \combweight_i = 1$ from \eqref{eq:appendix-jac}.

% \paragraph{Gradient.}
The shift gradient of $\jointdual$ is identical to \eqref{eq:appendix-centroid-grad}, while for the weights the chain rule gives
\begin{equation}
  \partial_{\weight_i} \jointdual
    = (\partial_{\weight_i} \combweight_i) \, (\partial_{\combweight_i} \energy) + \pvol_i
    = (\partial_{\combweight_i} \energy) + \pvol_i
    = -\bigl( \cellvol_i - \pvol_i \bigr) ,
  \label{eq:appendix-joint-grad}
\end{equation}
so that stacking the two gradients recovers the joint gradient \eqref{eq:joint-gradient}.

% \paragraph{Hessian.}
Differentiating a second time, the diagonal $(\shift_i,\shift_i)$ and $(\weight_i,\weight_i)$ blocks are exactly the Hessians we have already computed --- the centroid Hessian of $\centroiddual$ in \eqref{eq:centroid-hessian} and the volume Hessian of $\dualenergy$ in \eqref{eq:newton-step} --- but the joint problem also produces the mixed shift--weight blocks of \eqref{eq:joint-hessian}.
To handle all blocks at once, we introduced an auxliary variable that stacks the two multipliers of a site into $\mult_i = [\shift_i, \weight_i] \in \R^{d+1}$ and then define the affine field
\begin{equation}
  \psi_i(\xx) = 2\,\shift_i\tr (\xx - \tcentroid_i) + \weight_i = \mult_i\tr \phi_i(\xx) ,
  \qquad
  \phi_i(\xx) = \begin{bmatrix} 2(\xx - \tcentroid_i) \\ 1 \end{bmatrix} ,
  \label{eq:appendix-potential}
\end{equation}
through which $\mult_i$ enters the shifted power distance $\norm{\xx - \qq_i}^2 - \psi_i(\xx)$.
The Hessian is then the facet-based Laplacian assembled from $\phi_i$:
\begin{subequations}
  \label{eq:appendix-block}
  \begin{empheq}[left={\empheqlbrace\,}]{align}
    \partial^2_{\mult_i \mult_j} \jointdual
      &= \frac{1}{2}\frac{1}{\norm{\dual_i - \dual_j}} \int_{\face_{ij}} \phi_i(\xx)\, \phi_j(\xx)\tr \, dA,
    \label{eq:appendix-block-off}
    \\[2mm]
    \partial^2_{\mult_i \mult_i} \jointdual
      &= -\!\sum\nolimits_{j \sim i} \frac{1}{2}\frac{1}{\norm{\dual_i - \dual_j}} \int_{\face_{ij}} \phi_i(\xx)\, \phi_i(\xx)\tr \, dA.
    \label{eq:appendix-block-diag}
  \end{empheq}
\end{subequations}
Expanding $\phi_i$ splits \eqref{eq:appendix-block} into shift--shift, shift--weight, and weight--weight coefficients, with off-diagonal blocks
\begin{subequations}
  \label{eq:appendix-blocks-off}
  \begin{empheq}[left={\empheqlbrace\,}]{align}
    \partial^2_{\shift_i \shift_j\tr} \jointdual
      &= \frac{2}{\norm{\dual_i - \dual_j}} \int_{\face_{ij}}\! (\xx - \tcentroid_i)(\xx - \tcentroid_j)\tr \, dA , \label{eq:appendix-blocks-uu} \\
    \partial^2_{\shift_i \weight_j} \jointdual
      &= \frac{1}{\norm{\dual_i - \dual_j}} \int_{\face_{ij}}\! (\xx - \tcentroid_i) \, dA , \label{eq:appendix-blocks-uw} \\
    \partial^2_{\weight_i \weight_j} \jointdual
      &= \frac{1}{2}\frac{\area_{ij}}{\norm{\dual_i - \dual_j}} , \label{eq:appendix-blocks-ww}
  \end{empheq}
\end{subequations}
and diagonal blocks
\begin{subequations}
  \label{eq:appendix-blocks-diag}
  \begin{empheq}[left={\empheqlbrace\,}]{align}
    \partial^2_{\shift_i \shift_i\tr} \jointdual
      &= -\!\sum\nolimits_{j \sim i} \frac{2}{\norm{\dual_i - \dual_j}} \int_{\face_{ij}}\! (\xx - \tcentroid_i)(\xx - \tcentroid_i)\tr \, dA , \label{eq:appendix-blocks-uu-diag} \\
    \partial^2_{\shift_i \weight_i} \jointdual
      &= -\!\sum\nolimits_{j \sim i} \frac{1}{\norm{\dual_i - \dual_j}} \int_{\face_{ij}}\! (\xx - \tcentroid_i) \, dA , \label{eq:appendix-blocks-uw-diag} \\
    \partial^2_{\weight_i \weight_i} \jointdual
      &= -\!\sum\nolimits_{j \sim i} \frac{1}{2}\frac{\area_{ij}}{\norm{\dual_i - \dual_j}} , \label{eq:appendix-blocks-ww-diag}
  \end{empheq}
\end{subequations}
whose shift--shift, shift--weight, and weight--weight parts recover, respectively, the centroid Hessian \eqref{eq:centroid-hessian}, the mixed block \eqref{eq:joint-hessian}, and the volume Hessian \eqref{eq:newton-step}.

% \paragraph{Negative semidefiniteness and null space.}
Contracting the Hessian with a perturbation $\smash{\delta\mult = \{ (\delta\shift_i, \delta\weight_i) \}}$, again collapses each facet to a scalar square of the form
\begin{equation*}
  \delta\mult\tr \nabla^2 \jointdual \, \delta\mult
    = -\sum\nolimits_{i,j} \frac{1}{2\,\norm{\dual_i - \dual_j}} \int_{\face_{ij}} \bigl( \delta\psi_i(\xx) - \delta\psi_j(\xx) \bigr)^2 \, dA
    \;\le\; 0 ,
  \label{eq:appendix-joint-nsd}
\end{equation*}
now with variations of the affine field $\smash{\delta\psi_i(\xx) = 2\,\delta\shift_i\tr (\xx - \tcentroid_i) + \delta\weight_i}$.
Hence, $\nabla^2 \jointdual \preceq 0$ and $\jointdual$ is jointly concave.
As in Appendix~\ref{sec:appendix-centroid}, this quadratic form vanishes exactly when $\smash{\delta\psi_i - \delta\psi_j = \kappa_{ij}\,(a_i - a_j)}$ on every facet, now with $\smash{\delta a_i = \delta\psi_i}$ since the weights also vary.
The dilation and translation perturbation $\smash{\delta a_i = \mu\,a_i + 2\,\translation\tr\xx + c}$ always satisfies this condition and produces the only solution.
Matching the linear and constant parts as in \eqref{eq:appendix-centroid-kernel}, but now solving the constant part for the free weight perturbation, gives
\begin{equation*}
  \delta\shift_i = \mu\,\dual_i + \translation ,
  \qquad
  \delta\weight_i = \mu\,\gamma_i + 2\,\translation\tr\tcentroid_i + c ,
\end{equation*}
which is exactly the $(d\!+\!2)$-dimensional family in \eqref{eq:joint-nullspace}.
One can verify directly that it lies in the null-space, since $\smash{\delta\psi_i - \delta\psi_j = \mu\,(a_i - a_j)}$ vanishes on $\face_{ij}$.
There, $\mu$ scales the lifted shifted sites $\smash{(\dual_i, \sheight_i)}$, $\translation$ translates them horizontally, and $c$ translates them vertically, which together form the homotheties \eqref{eq:infinitesimal-homothety} that leave the power diagram unchanged (\S\ref{sec:background}).
In contrast to $\centroiddual$, these modes are always present because the free weights absorb the weight changes they require.

\section{Existence of solutions}
\label{sec:appendix-existence}

On the open set of distinct shifted sites with non-empty cells, the energies $\centroiddual$ and $\jointdual$ are differentiable and concave, and their stationary points are exactly the power diagrams that meet the constraints.
We detail here which targets admit such a stationary point, for each formulation.

\paragraph{Volume constraints.}
The $n$ weights of $\dualenergy$, up to a common constant, match the $n$ volume constraints, up to their sum, and any positive volumes summing to $\domainVol$ are realized for any distinct sites \cite{Aurenhammer1998Minkowski}.

\paragraph{Centroid constraints.}
The $nd$ shifts of $\centroiddual$ likewise match the $nd$ centroid constraints, since the free volumes leave these constraints independent.
When the pivots are not coplanar, the Hessian of $\centroiddual$ is negative definite (\S\ref{sec:centroid-constraints}), so by the inverse function theorem, solutions persist under small perturbations of the targets.
Not every target is realizable, however.
In $\domain = [0,1]$ with targets $\smash{\tcentroid = (0.3, 0.7)}$, the partition $[0.05, 0.55]$ and its complement realizes the targets, but no power diagram does, since its interval cells $[0,s]$ and $[s,1]$ would require both $s = 0.6$ and $s = 0.4$; the supremum of $\centroiddual$ is then approached as the two shifted sites merge.

\paragraph{Volume and centroid constraints.}
Since $\jointdual$ is the minimum over partitions of the Lagrangian \eqref{eq:joint-lagrangian}, which is affine in the multipliers for fixed cells, it varies linearly along the modes \eqref{eq:joint-nullspace}, which leave the cells unchanged.
A maximizer therefore requires $\jointdual$ to be bounded along these modes.
Along the constant and translation modes, the slope depends only on the data and vanishes exactly under the identities \eqref{eq:partition-feasibility}.
Along the dilation mode, which scales the shifted sites by a factor $s \!>\! 0$, the slope instead depends on the configuration: once the volumes are enforced, it equals $\smash{2 \sum_i \pvol_i (\tcentroid_i - \centroid_i)\tr \dual_i}$ \cite[Theorem~4.2]{Bourne2025Inverting}.
Maximizing $\jointdual$ over the weights gives $2\helper(\dual)$ up to a constant, with $\smash{\helper(\dual) = \sum_i \pvol_i\,\dual_i\tr \tcentroid_i - \sigma_{\achievable}(\dual)}$, where $\smash{\sigma_{\achievable}(\dual) = \max_{\centroid \in \achievable} \sum_i \pvol_i\,\dual_i\tr \centroid_i}$ is the support function of $\achievable$ \cite[Prop.~3.5]{Bourne2026ConvexOrder}, so $\helper$ is positively $1$-homogeneous in the shifted sites.
Its supremum is therefore $+\infty$ if $\tcentroid \notin \achievable$, approached as the shifted sites dilate along a separating direction, and $0$ otherwise.
This value is attained at distinct shifted sites exactly when $\dual$ is a normal of $\achievable$ at $\tcentroid$ with distinct sites, \ie when $\tcentroid$ is an exposed point of $\achievable$ (\S\ref{sec:background}); elsewhere in $\achievable$, it is only approached as the shifted sites merge.
This additional condition requires the targets to be exposed points of $\achievable$, which lie on its relative boundary, a hypersurface of the affine subspace prescribed by \eqref{eq:partition-feasibility}, and thus fixes the overall spread of the target centroids relative to the prescribed volumes, but not the shape of their configuration.
In the unit interval with two cells of volume $\smash{\tfrac12}$, for instance, the targets $(0.3, 0.7)$ satisfy \eqref{eq:partition-feasibility}, but must be spread by a factor $\smash{\tfrac54}$ about $\smash{\tfrac12}$, to $\smash{(\tfrac14, \tfrac34)}$, to be realized \cite[Example~4.7]{Bourne2025Inverting}.

\section{Softened centroid constraints}
\label{sec:appendix-regularized}

\paragraph{Duality.}
We introduce the shifts $\shift$ and weights $\weight$ as multipliers of the constraints \eqref{eq:regularized-primal-cons} and \eqref{eq:regularized-primal-vol}, scaling the shifts by two as in \S\ref{sec:centroid-constraints}.
The Lagrangian of \eqref{eq:regularized-primal} is then the joint Lagrangian \eqref{eq:joint-lagrangian} augmented by the slack terms,
\begin{equation*}
  \lag_\alpha(\cell, \slack, \shift, \weight)
    = \lag(\cell, \shift, \weight) + \sum\nolimits_i \pvol_i \bigl( \alpha \norm{\slack_i}^2 + 2\,\shift_i\tr \slack_i \bigr) .
\end{equation*}
The slack terms are decoupled from the partition and strictly convex in each $\slack_i$ for $\alpha \!>\! 0$, so their minimum is attained at $\smash{\slack_i = -\shift_i / \alpha}$ with value $\smash{-\tfrac{1}{\alpha} \pvol_i \norm{\shift_i}^2}$.
Relaxing the partitions to fractional ones, \ie to cell indicator functions $\basis_i(\xx) \in [0,1]$ with $\sum_i \basis_i = 1$, the Lagrangian is linear in the partition and concave in the multipliers, so the minimization and maximization can be exchanged.
Minimizing over partitions first then yields $\jointdual$ of \eqref{eq:joint-dual}, and adding the slack minimum gives $\jointdual_\alpha$ in \eqref{eq:regularized-dual}.
Its derivatives follow from those of $\jointdual$ in Appendix~\ref{sec:appendix-joint}, with the extra term $\smash{-\tfrac{2}{\alpha}\pvol_i\,\shift_i}$ in the shift gradient and $\smash{-\tfrac{2}{\alpha}\pvol_i\,\Id}$ in the diagonal shift--shift blocks.

\paragraph{Existence and bounds.}
Let $\cell^\star$ be any partition with volumes $\smash{\pvol_i}$, and denote its centroid moments by $\smash{\mathbf{m}_i^\star = \int_{\cell^\star_i} (\xx - \tcentroid_i) \, d\xx}$.
Since $\jointdual(\shift, \weight)$ is a minimum over partitions, evaluating the Lagrangian at $\cell^\star$ gives $\smash{\jointdual(\shift, \weight) \le C - 2\sum_i \shift_i\tr \mathbf{m}_i^\star}$ for a constant $C$ independent of the multipliers, so $\jointdual_{\alpha} \to -\infty$ as $\norm{\shift} \to \infty$, uniformly in the weights.
For bounded shifts, $\jointdual_{\alpha}$ is, up to the additive constant on the weights, coercive in the weights whenever $\smash{\sum_i \pvol_i = \domainVol}$, as in the volume-constrained problem \eqref{eq:volume-dual}.
Hence $\jointdual_{\alpha}$ attains its maximum for any targets.
Since $\jointdual_{\alpha}$ is strongly concave in the shifts, two maximizers $\smash{(\shift^1, \weight^1)}$ and $\smash{(\shift^2, \weight^2)}$ would give at their midpoint a value exceeding the maximum by at least $\smash{\tfrac{1}{4\alpha}\sum_i \pvol_i \norm{\shift^1_i - \shift^2_i}^2}$, so the maximizing shifts are unique.
At a maximizer, optimality in the weights enforces the volumes $\smash{\pvol_i}$, also for fractional partitions, and optimality in the shifts gives $\smash{\shift_i = \alpha\,(\tcentroid_i - \centroid_i)}$, where $\centroid_i$ is the centroid of the possibly fractional cell $i$.
For targets in the convex hull of $\domain$, both centroids lie in this hull, so $\smash{\norm{\shift_i} \le \alpha \operatorname{diam}(\domain)}$.
Therefore, $\smash{\norm{\dual_i - \dual_j} \ge \norm{\qq_i - \qq_j} - \norm{\shift_i} - \norm{\shift_j} > 0}$ whenever \eqref{eq:regularized-distinct} holds, in which case the optimal partition is the power diagram of the shifted weighted sites.

\paragraph{Line-constrained maximization.}
The constrained integrator of \S\ref{sec:flow-constrained} maximizes $\jointdual$ subject to the linear constraints \eqref{eq:flow-line-constraint}.
We place the origin at $\domainCentroid$, so that $\smash{\sum_i \pvol_i\,\centroid_i = 0}$ on $\achievable$, and we consider a direction $\mathbf{d}$ with $\smash{\sum_i \pvol_i\,\mathbf{d}_i = 0}$, sites $\qq$ with $\smash{\kappa = \sum_i \pvol_i\,\mathbf{d}_i\tr \qq_i > 0}$ and $\smash{\mathbf{m} = \sum_i \pvol_i\,\qq_i}$, and arbitrary targets $\pred$, with $\smash{\bar{\translation} = \sum_i \pvol_i\,\pred_i / \domainVol}$.

\begin{lemma}
  \label{lem:line-constrained}
  Consider the maximization of $\jointdual$ with sites $\qq$ and targets $\pred$ subject to $\smash{\sum_i \pvol_i\,\mathbf{d}_i\tr \shift_i = 0}$ and $\smash{\sum_i \pvol_i\,\shift_i = 0}$, with multipliers $s$ and $\translation$ moving the targets to $\smash{\pred - s\,\mathbf{d} - \translation}$.
  Its dual is the minimization of $s$ over $(s, \translation)$ subject to $\smash{\pred - s\,\mathbf{d} - \translation \in \achievable}$, whose feasible translation is always $\smash{\translation = \bar{\translation}}$.
\end{lemma}
\begin{proof}
  As in Appendix~\ref{sec:appendix-existence}, maximizing over the weights gives $\smash{2\,\helper(\dual) + \text{const}}$, now with $\smash{\helper(\dual) = \sum_i \pvol_i\,\dual_i\tr \pred_i - \sigma_{\achievable}(\dual)}$, and the constraints read $\smash{\sum_i \pvol_i\,\mathbf{d}_i\tr \dual_i = \kappa}$ and $\smash{\sum_i \pvol_i\,\dual_i = \mathbf{m}}$ in terms of the generators $\dual = \qq + \shift$.
  Dualizing the constraints with multipliers $2s$ and $2\translation$ gives $\smash{\max_{\dual} 2 \bigl( \sum_i \pvol_i\,\dual_i\tr (\pred_i - s\,\mathbf{d}_i - \translation) - \sigma_{\achievable}(\dual) \bigr) + 2s\kappa + 2\translation\tr\mathbf{m}}$.
  Since $\sigma_{\achievable}$ is the support function of $\achievable$, the inner maximum is $0$ if $\smash{\pred - s\,\mathbf{d} - \translation \in \achievable}$ and $+\infty$ otherwise.
  Every point of $\achievable$ verifies $\smash{\sum_i \pvol_i\,\centroid_i = 0}$, and $\mathbf{d}$ is centered, so this condition forces $\smash{\translation = \bar{\translation}}$ for every $s$, and the dual minimizes $\smash{2s\kappa + 2\bar{\translation}\tr\mathbf{m}}$ over the feasible $s$, and $\kappa \!>\! 0$ gives the claim.
\end{proof}
In terms of partitions, this dual reads
\begin{subequations}
  \label{eq:flow-line-primal}
  \begin{alignat}{2}
    & \min_{\cell, s, \translation} \quad && \sum\nolimits_i \int_{\cell_i} \norm{\xx - \qq_i}^2 \, d\xx \label{eq:flow-line-primal-obj}
  \end{alignat}
  \vspace{-\belowdisplayskip}\vspace{-\abovedisplayskip}
  \begin{empheq}[left={\text{s.t.} \quad \empheqlbrace}]{alignat=1}
    & \int_{\cell_i} \bigl( \xx - \pred_i + s\,\mathbf{d}_i + \translation \bigr) \, d\xx = 0 \ \ \forall i , \label{eq:flow-line-primal-cons} \\
    & \cellvol_i = \pvol_i \ \ \forall i , \label{eq:flow-line-primal-vol}
  \end{empheq}
\end{subequations}
since the Lagrangian of \eqref{eq:flow-line-primal} is linear in $s$ and $\translation$, and is bounded below in them exactly when the shifts verify the two linear constraints.
Under these constraints, the objective \eqref{eq:flow-line-primal-obj} is affine in the centroids and equals $\smash{2s\kappa}$ up to a constant, while summing \eqref{eq:flow-line-primal-cons} over the cells fixes the translation independently of $s$.
In \S\ref{sec:flow-constrained}, the sites are $\qq^k$, the targets are $\smash{\pred^{k+1}}$, and the direction is $\dual^k$; if $\dual^k$ is not centered, the second constraint removes its mean $\smash{\bar\dual^k = \sum_i \pvol_i\,\dual^k_i / \domainVol}$ from the first one, so the lemma applies with $\smash{\mathbf{d} = \dual^k - \bar\dual^k}$, the translation absorbing $\smash{s\,\bar\dual^k}$.
Finally, $\kappa \!>\! 0$, since $\qq^k$ maximizes $\smash{\sum_i \pvol_i\,\mathbf{d}_i\tr \centroid_i}$ over $\achievable$, which contains the centroids $\smash{\varepsilon\,\mathbf{d}}$ of a fractional partition with volumes $\pvol$ for $\varepsilon \!>\! 0$ small enough, at which this function is positive since the generators are pairwise distinct.

\section{Feasible volumes for prescribed centroids}
\label{sec:appendix-maxent}

In the centroid-only problem of \S\ref{sec:centroid-constraints}, the volumes are unknown, and the conservation identities \eqref{eq:partition-feasibility} are satisfiable only if there exist non-negative volumes $\cellvol_i$ with $\smash{\sum_i \cellvol_i \tcentroid_i = \domainVol\,\centroid_\domain}$ and $\smash{\sum_i \cellvol_i = \domainVol}$.
Drawing inspiration from \citet{Hormann2008MaxEnt}, one way to recover such volumes is to solve
\begin{subequations}
  \label{eq:maxent-feasibility}
  \begin{alignat}{2}
    & \min_{\{\cellvol_i\}} \quad && \sum\nolimits_i \cellvol_i \, \log\bigl( \cellvol_i / \domainVol \bigr) \label{eq:maxent-feasibility-obj} \\
    & \text{s.t.} \quad && \sum\nolimits_i \cellvol_i \, \tcentroid_i = \domainVol \, \centroid_\domain , \quad \sum\nolimits_i \cellvol_i = \domainVol , \label{eq:maxent-feasibility-cons}
  \end{alignat}
\end{subequations}
which is equivalent to the \emph{unconstrained} convex minimization
\begin{equation}
  \min_{\yy \in \R^{d+1}} \
    \domainVol \sum\nolimits_i \exp\bigl( -\yy\tr \lift_i \bigr)
    + \yy\tr \lift_\domain ,
  \label{eq:maxent-dual}
\end{equation}
with the homogeneous lifts $\lift_i \!=\! [\tcentroid_i, 1]$ and $\lift_\domain \!=\! \domainVol\,[\centroid_\domain, 1]$, whose minimizer recovers the volumes as $\smash{\cellvol_i = \domainVol \exp(-\yy\tr \lift_i)}$.

\section{Bordered Newton step}
\label{sec:appendix-bordered}

We detail the Newton step of the constrained integrator (\S\ref{sec:flow-constrained}), which maximizes \eqref{eq:flow-line-constraint}.
Denoting the negated Hessian blocks of $\jointdual$ by $\smash{\Kmat = -\partial^2_{\shift\shift}\jointdual}$ and $\smash{\wlap = -\partial^2_{\weight\weight}\jointdual}$, the mixed block \eqref{eq:joint-hessian} by $\smash{\Bmat = \partial^2_{\shift\weight}\jointdual = -\partial\cellvol/\partial\shift}$, and the constraint rows by $\smash{\Cmat\,\shift = \bigl( \sum_i \pvol_i\,\shift_i\tr \dual^k_i, \, \sum_i \pvol_i\,\shift_i \bigr)}$, this bordered system reads
\begin{equation}
  \begin{bmatrix} \Kmat & -\Bmat & \Cmat\tr \\ -\Bmat\tr & \wlap & 0 \\ \Cmat & 0 & 0 \end{bmatrix}
  \begin{bmatrix} \delta\shift \\ \delta\weight \\ (2s, 2\translation) \end{bmatrix}
  =
  \begin{bmatrix} \partial_{\shift}\jointdual \\ \partial_{\weight}\jointdual \\ 0 \end{bmatrix} ,
  \label{eq:flow-constrained-newton}
\end{equation}
where the last row vanishes since every iterate satisfies the linear constraints, and where we pin one weight to fix the additive constant.
At a stationary point, the updates vanish and the first row reduces to $\smash{\partial_{\shift}\jointdual = \Cmat\tr (2s, 2\translation)}$, \ie to the correction \eqref{eq:flow-shake-step-corr} with the new particles as centroids.

We solve this system by block elimination of the multipliers.
Since the negated Hessian of $\jointdual$ is singular along the homotheties \eqref{eq:joint-nullspace}, we first add to it the outer product of the constraint rows $\smash{\hat\Cmat = [\Cmat \;\, 0]}$, padded with zeros on the weights, which yields
\begin{equation}
  \Nmat = \begin{bmatrix} \Kmat & -\Bmat \\ -\Bmat\tr & \wlap \end{bmatrix} + \hat\Cmat\tr \hat\Cmat .
  \label{eq:appendix-bordered-matrix}
\end{equation}
This modification leaves the solution unchanged, since the updates satisfy $\smash{\hat\Cmat\,(\delta\shift, \delta\weight) = 0}$, so the added term vanishes on them.
With one weight pinned, the only remaining null-space modes of the negated Hessian are the dilation and the translations of \eqref{eq:joint-nullspace}, on which $\smash{\hat\Cmat}$ evaluates to the $(d+1) \times (d+1)$ matrix with entries $\smash{\sum_i \pvol_i\,(\dual^k_i)\tr \dual_i}$, $\smash{\sum_i \pvol_i\,\dual^k_i}$, $\smash{\sum_i \pvol_i\,\dual_i}$ and $\smash{\domainVol\,\Id}$, where $\dual$ are the current generators.
Its determinant is proportional to $\smash{\sum_i \pvol_i\,(\dual^k_i)\tr (\dual_i - \bar\dual)}$, with $\smash{\bar\dual = \sum_i \pvol_i\,\dual_i / \domainVol}$, so $\Nmat$ is nonsingular whenever this correlation is nonzero, as in the condition $\kappa \!>\! 0$ of Lemma~\ref{lem:line-constrained}.

Denoting by $\residual = (\partial_{\shift}\jointdual, \partial_{\weight}\jointdual)$ the gradient, the Newton step then follows in three stages,
\begin{subequations}
  \label{eq:flow-constrained-elimination}
  \begin{empheq}[left={\empheqlbrace\,}]{align}
    \Nmat \, \bigl[ \boldsymbol{\xi} , \, \Xi \bigr] &= \bigl[ \residual , \, \hat\Cmat\tr \bigr] , \label{eq:flow-constrained-elimination-solve} \\[1mm]
    \bigl( \hat\Cmat\, \Xi \bigr) \, (2s, 2\translation) &= \hat\Cmat\, \boldsymbol{\xi} , \label{eq:flow-constrained-elimination-gauge} \\[1mm]
    (\delta\shift, \delta\weight) &= \boldsymbol{\xi} - \Xi\, (2s, 2\translation) . \label{eq:flow-constrained-elimination-update}
  \end{empheq}
\end{subequations}
The first stage solves $\Nmat$ for $d + 2$ right-hand sides: the gradient $\residual$, and the columns of $\smash{\hat\Cmat\tr}$, which are $\smash{\pvol_i\,\dual^k_i}$ for the dilation and $\smash{\pvol_i\,\Id}$ for the translations on the shifts.
The second stage solves a dense $(d + 1) \times (d + 1)$ symmetric positive-definite system for the gauge $(s, \translation)$, \ie the negated Schur complement of the bordered system with $\Nmat$ in place of the negated Hessian, and the last stage recovers the updates of the shifts and weights.
Since $\smash{\hat\Cmat\tr \hat\Cmat}$ is dense with rank $d + 1$, $\Nmat$ is a sparse matrix plus a low-rank term, which can be applied without being formed, \eg within conjugate gradients.
% TODO: choose the solver for N to match ../src (currently only a sparse Cholesky of H + eps M); a sparse direct factorization needs the dense rank-(d+1) term handled separately.

\end{document}
